\section{Background}
\label{sec:background}

\subsection{GraphCast}

GraphCast~\cite{graphcast} is an encode--process--decode graph network over the
global atmospheric state. The encoder maps the input fields (atmospheric variables
on pressure levels plus surface fields, on a latitude--longitude grid) onto a
multi-scale icosahedral mesh. A 16-layer message-passing processor, with 512 latent
dimensions per node, propagates information along mesh edges of varying length;
this is where the learned dynamics live. The decoder maps mesh states back to the
grid as a 6-hour forecast, and longer forecasts come from feeding the model its own
output.

Two architectural facts matter later. Rolling the model forward (or teacher-forcing
it through consecutive observed states) supplies a time axis over its internal
features, which the causal-discovery route depends on (Sec.~\ref{sec:temporal}).
And the multi-mesh is heterogeneous: refinement leaves some nodes much better
connected than others. That heterogeneity is a plausible origin for the
position-coding (``grid-locked'') features reported in~\cite{macmillan2025}, and it
is why the synthetic forecaster we build in Sec.~\ref{sec:gnn} uses a heterogeneous
mesh as well.

\subsection{Sparse autoencoders}

An SAE is an overcomplete dictionary trained to reconstruct a layer's activations
from a sparse code, in the hope that the dictionary directions are monosemantic
where the raw neurons are not~\cite{cunningham2023,bricken2023}. The TopK variant
keeps the $K$ largest feature activations per sample. \citet{macmillan2025}
obtained their GraphCast features this way, from processor activations. Their
evaluation stops at feature existence and single-feature interventions; ours
continues to the causal graph among the features.

\subsection{Causal discovery from time series}

We use three discovery methods whose assumptions barely overlap.
PCMCI~\cite{runge2019} is the primary one: a condition-selection stage followed by
momentary-conditional-independence tests, designed for autocorrelated
high-dimensional time series, returning lagged directed edges with significance
levels. PCMCI orients only lagged edges ($\tau \ge 1$); its successor
PCMCI+~\cite{runge2020} also estimates contemporaneous ($\tau = 0$) adjacencies,
which becomes necessary once dynamics act faster than the sampling interval
(Sec.~\ref{sec:subsample}). TSCI~\cite{tsci} comes from the
convergent-cross-mapping tradition and infers causation from reconstructed
attractor geometry. DYNOTEARS~\cite{dynotears} learns a weighted dynamic Bayesian
network by acyclicity-constrained continuous optimization. \note{DYNOTEARS runs in
an isolated venv here due to a numpy/pandas pin; see \texttt{requirements.txt}.}
Where the three agree we gain confidence; where they disagree, the disagreement
itself tells us which assumptions bind.

One identifiability fact shapes several experiments below. Under linear dynamics
with Gaussian noise, the direction of a contemporaneous edge cannot be determined
from observational data: $X \to Y$ and $Y \to X$ induce the same covariance.
Shimizu et al.~\cite{shimizu2006} showed that non-Gaussian noise breaks the tie
(the LiNGAM result: linear, non-Gaussian, acyclic models are fully identifiable,
because higher-order moments carry directional information). The standard ParCorr
test inside PCMCI+ uses second-order statistics only and so cannot exploit this;
Sec.~\ref{sec:subsample} measures what that costs.

\subsection{Parameter decomposition}
\label{sec:pd-background}

SAEs analyze activations. For the weights we use stochastic parameter
decomposition (SPD)~\cite{spd}, which writes each weight matrix as a sum of
rank-one components, $W \approx \sum_c U_{:,c} V_{c,:}$, trained so that the
components sum to $W$, few are needed for any given input, and each is simple
(rank one, single layer). ``Needed'' has an operational definition here. A gate
$g_c(x) \in [0,1]$, learned by randomly masking components and requiring the
masked network to reproduce the original output, estimates how far component $c$
can be attenuated at input $x$ before the computation changes. VPD~\cite{vpd}
hardens the scheme with adversarially chosen masks, so faithfulness must hold
under worst-case ablations, and adds a frequency penalty that pushes each
component to be needed rarely---that is, to specialize.

The gate is the reason we adopted this method. It is an ablation, so causal
importance comes built in; no separate intervention experiment is needed to ask
whether a component is used. Better still for our purposes, when the decomposed
layer is applied at every spatial position with shared weights, as in message
passing, evaluating the gate per position gives each component a spatial usage
map---and the model contains no per-position parameters that could fake this
(Sec.~\ref{sec:vpd}).

\subsection{SAVAR}

SAVAR (Spatially Aggregated VAR)~\cite{savar} is a stochastic spatiotemporal
climate benchmark: a handful of spatial \emph{modes} evolve under a
vector-autoregressive process with a prescribed causal graph and are rendered onto
a grid through fixed mode weight maps. The graph among modes is known by
construction, which is what lets us score a discovery pipeline before trusting it
on GraphCast.

Our default instantiation, all of it part of the given data-generating process:
\begin{itemize}[nosep]
  \item $N=8$ latent modes on a $50\times50$ grid ($L=2500$ pixels).
  \item Mode maps $W$: eight fixed, $L_1$-normalized, non-overlapping Gaussian
        blobs in a $3\times3$ layout. $W$ maps modes to pixels;
        $W^{+}=\mathrm{pinv}(W)$ maps pixels back to modes.
  \item The causal graph $G$: a hand-designed coefficient set
        $\{X_i(t-\tau)\!\to\!X_j(t)\}$, each mode carrying an autoregressive term
        plus cross-edges. It was engineered to contain a hub (X0), a converging
        node with two parents (X3), a pure sink (X7), a weak lag-2 feedback
        (X2$\to$X0), negative edges (X0$\to$X5, X2$\to$X3, X3$\to$X7), and a mix of
        lag-1 and lag-2 structure.
  \item Dynamics and noise: $Z(t)=\sum_\tau \Phi(\tau)Z(t-\tau)+\xi(t)$,
        observations $=W^{+}Z+\text{spatial noise}$, with $\lambda=1$, $D_x=I_N$,
        $D_y=I_L$, giving $\Sigma_y=W^{+}W^{+\top}+I_L$; $T=500$ usable steps after
        a burn-in of $200$. The 100 realisations share $W$ and $G$ and differ only
        in noise seed.
\end{itemize}
The coefficient dictionary $G$ is the ground truth every method in
Sec.~\ref{sec:methodcompare} is scored against. This default generator is linear
and Gaussian; Sec.~\ref{sec:ext} makes it harder in targeted ways (fine temporal
cadence with heterogeneous lags, nonlinear dynamics, skewed innovations, widely
spread mode timescales), each aimed at a question the default cannot pose.
